\chapter{Evaluation, Testing, and MLOps}
\label{ch:evaluation-mlops}

\begin{learningobjectives}
After this chapter, you should be able to:
\begin{itemize}
    \item translate a real decision into a measurable evaluation contract;
    \item choose leakage-resistant data splits and task-appropriate metrics;
    \item quantify sampling uncertainty instead of reporting a lonely score;
    \item test data, numerical code, learned behaviour, and the surrounding
          system;
    \item design an experiment in which a stronger model can genuinely lose;
    \item define release gates and production monitoring for an ML service.
\end{itemize}
\end{learningobjectives}

\begin{plainlanguage}
Training asks, ``Can the optimiser reduce this loss on these examples?''
Evaluation asks, ``Does the resulting system help on the cases that matter,
under conditions that resemble use, without unacceptable side effects?''
Those are different questions.
The second one is the engineering question.
\end{plainlanguage}

\section{Begin with an Evaluation Contract}

Evaluation begins before a model is selected.
A useful \keyterm{evaluation contract} records the intended decision, the unit
being predicted, the population of interest, the costs of errors, and the
conditions under which the result will be used.

\begin{tcbitemize}[skin=sectionraster]
    \tcbitem[title=Decision and action, raster multicolumn=3]
    What action changes when the prediction changes?
    A probability estimate that never changes a decision may be analytically
    interesting but operationally irrelevant.

    \tcbitem[title=Unit and population, raster multicolumn=3]
    Is one row a machine cycle, a patient, an invoice, a customer, or a token?
    Define who or what must generalise beyond the collected sample.

    \tcbitem[title=Error costs, raster multicolumn=3]
    A false alarm and a missed event rarely cost the same amount.
    Record monetary, safety, fairness, and workload consequences.

    \tcbitem[title=Operating conditions, raster multicolumn=3]
    State expected input ranges, latency and memory budgets, human review,
    fallback behaviour, and prohibited uses.
\end{tcbitemize}

\begin{engineernote}
Write the acceptance criteria before training.
If the criteria are invented after looking at results, model selection becomes
an exercise in moving the goalposts with impressive numerical precision.
\end{engineernote}

\subsection{Baseline ladder}

Every candidate should be compared with a ladder of increasingly credible
baselines:

\begin{enumerate}
    \item a random or majority-class rule that exposes broken metrics;
    \item a domain heuristic that represents current practice;
    \item a simple statistical model that is easy to inspect;
    \item the current production or previously accepted model;
    \item the proposed model under exactly the same evaluation protocol.
\end{enumerate}

A complex candidate is not entitled to win.
If logistic regression is more accurate, better calibrated, faster, and easier
to maintain, the engineering victory belongs to logistic regression.

\section{Data Splits and Leakage}

The usual train/validation/test vocabulary describes different permissions.
The \keyterm{training set} may influence learned parameters.
The \keyterm{validation set} may influence hyperparameters, thresholds, model
choice, and stopping.
The \keyterm{test set} may influence the final claim, but not the system that
is being claimed about.

\begin{table}[htbp]
\centering
\caption{Choose a split that imitates the intended generalisation boundary.}
\label{tab:split-strategies}
\begin{tblr}{
    width=\textwidth,
    colspec={Q[l,0.18\textwidth] X[l] X[l]},
    row{1}={font=\bfseries, bg=\maincolor!15!white},
    hlines,
    vlines}
Split & Use when & Typical leakage prevented \\
Random stratified & Independent and identically distributed rows &
Class proportions drifting between splits \\
Grouped & Several rows belong to one person, machine, site, or source &
The same entity appearing in train and test \\
Temporal & The system predicts future observations &
Future information influencing the past \\
Geographic or site holdout & Deployment expands to unseen locations &
Site-specific shortcuts and duplicated procedures \\
Nested cross-validation & Hyperparameter selection and performance estimation
share scarce data &
Optimism from selecting and reporting on the same folds \\
Private challenge set & Repeated model development risks benchmark overfitting &
Direct and indirect adaptation to known test cases \\
\end{tblr}
\end{table}

\keyterm{Data leakage} occurs when information unavailable at the intended
prediction time influences training, preprocessing, selection, or evaluation.
It can be obvious, such as leaving the target in the features, or subtle, such
as fitting a scaler on all rows before cross-validation.

\begin{cautionbox}[The pipeline must be inside the split]
Imputation, normalisation, vocabulary construction, feature selection,
oversampling, and dimensionality reduction learn from data.
Fit those operations on each training fold and apply the fitted transformation
to its validation fold.
Libraries such as scikit-learn provide pipelines precisely to preserve this
boundary.\textsuperscript{\cite{pedregosaScikitlearnMachineLearning2011}}
\end{cautionbox}

\subsection{Contamination and deduplication}

For web-scale and generative models, exact train--test separation may be hard
to prove.
Record dataset provenance, licenses, hashes, and collection dates; perform
exact and near-duplicate searches; and treat an unexplained benchmark jump as
a reason to investigate rather than a reason to celebrate.
Datasheets make these questions explicit at dataset level, while model cards
record intended use and evaluation
context.\textsuperscript{\cite{gebruDatasheetsDatasets2021,
mitchellModelCardsModel2019}}

\section{Metrics: Translate Errors into Numbers}

No metric is universally best.
A metric is a compressed statement about which distinctions and errors the
evaluation considers important.

\subsection{Classification}

For a binary classifier, let $TP$, $FP$, $TN$, and $FN$ denote true positives,
false positives, true negatives, and false negatives.

\begin{align}
    \text{precision} &= \frac{TP}{TP+FP},
    &\text{recall} &= \frac{TP}{TP+FN},
    \label{eq:precision-recall}\\
    \text{specificity} &= \frac{TN}{TN+FP},
    &F_1 &= 2\frac{\text{precision}\cdot\text{recall}}
                    {\text{precision}+\text{recall}}.
    \label{eq:specificity-f1}
\end{align}

Accuracy answers, ``What fraction was labelled correctly?''
It can be misleading when classes are imbalanced or error costs differ.
If only one in a thousand components fails, predicting ``healthy'' every time
achieves \qty{99.9}{\percent} accuracy and catches no failures at all.

\begin{workedexample}[A maintenance threshold]
Suppose a model raises 50 alarms.
Of these, 35 correspond to real failures, and it misses another 15 failures.
Then $TP=35$, $FP=15$, and $FN=15$.
Precision and recall are both $35/50=0.70$, so $F_1=0.70$.
Whether this is acceptable cannot be inferred from $F_1$ alone: the decision
depends on the cost of inspection, downtime, and a missed hazardous failure.
\end{workedexample}

Ranking metrics such as ROC AUC and average precision evaluate performance
across thresholds.
Operational decisions still require a threshold, so report results at the
intended operating point and include the resulting workload.

\subsection{Regression and forecasting}

For errors $e_i=y_i-\hat y_i$ over $n$ observations,

\begin{align}
    \operatorname{MAE} &= \frac{1}{n}\sum_{i=1}^{n}|e_i|,
    \label{eq:mae}\\
    \operatorname{MSE} &= \frac{1}{n}\sum_{i=1}^{n}e_i^2,
    &\operatorname{RMSE} &= \sqrt{\operatorname{MSE}}.
    \label{eq:mse-rmse}
\end{align}

MAE is in the target's units and weights absolute errors linearly.
RMSE is also in the target's units but penalises large errors more strongly.
The coefficient of determination,

\begin{equation}
    R^2 = 1 -
    \frac{\sum_i(y_i-\hat y_i)^2}
         {\sum_i(y_i-\bar y)^2},
    \label{eq:r-squared}
\end{equation}

compares squared error with a mean-prediction baseline.
$R^2$ can be negative on test data when the candidate is worse than that
baseline.
This is not a software error; it is a useful message from the data.

For forecasting, calculate errors with rolling or expanding temporal origins.
A random split allows the model to learn from a future that the deployed
system will never possess.

\subsection{Probabilities and calibration}

A probability is useful when statements such as ``about 20 out of 100 cases
assigned probability 0.2 occur'' are empirically true.
This agreement is \keyterm{calibration}.
Evaluate it with reliability diagrams and proper scoring rules such as log
loss or the Brier score, rather than checking only the winning class.

For a binary target $y_i\in\{0,1\}$ and predicted probability $p_i$,

\begin{equation}
    \operatorname{Brier} = \frac{1}{n}\sum_{i=1}^{n}(p_i-y_i)^2.
    \label{eq:brier-score}
\end{equation}

Calibration must be checked on data not used to fit the calibrator.
It should also be checked by relevant subgroup when error consequences differ.

\section{Uncertainty, Resampling, and Honest Comparisons}

A score calculated from a finite sample is a random quantity.
Reporting more decimal places does not create more evidence.

\subsection{Bootstrap interval}

Given a test set of $n$ independent evaluation units:

\begin{enumerate}
    \item sample $n$ units with replacement;
    \item calculate the metric on the resample;
    \item repeat many times using a recorded random seed;
    \item report suitable quantiles of the bootstrap distribution.
\end{enumerate}

Resample at the independent unit.
If one patient contributes many images, resampling individual images invents
more independent evidence than the experiment contains.

\begin{engineernote}
For a paired model comparison, resample cases once and calculate both models
on each same resample.
The distribution of score differences is usually more informative than two
separate confidence intervals.
\end{engineernote}

\subsection{Seeds and repeated experiments}

Neural training and reinforcement learning can vary substantially across
random initialisations, data order, and environment trajectories.
Report all planned seeds, the aggregation rule, and dispersion.
Do not quietly select the seed that tells the best story.

Hyperparameter search consumes validation evidence.
The final test set should be opened only after the model, preprocessing,
threshold, and stopping rule are frozen.
If the test result triggers another design iteration, it has become validation
data and a fresh test is required.

\section{A Testing Pyramid for Learned Systems}

Traditional software specifies behaviour in code.
ML systems distribute behaviour across code, data, parameters, configuration,
and infrastructure, which expands the test surface and creates hidden
technical debt.\textsuperscript{\cite{sculleyHiddenTechnicalDebt2015}}

\begin{center}
\begin{tikzpicture}[
    every node/.style={font=\sffamily\small,align=center},
    level/.style={draw=\maincolor!70!black,fill=\maincolor!8!white,
                  minimum height=9mm}]
    \node[level,minimum width=4.5cm] at (0,0)
        {Decision and safety validation};
    \node[level,minimum width=6.5cm] at (0,-0.95)
        {End-to-end and human-in-the-loop tests};
    \node[level,minimum width=8.5cm] at (0,-1.90)
        {Model behaviour and robustness tests};
    \node[level,minimum width=10.5cm] at (0,-2.85)
        {Data, numerical, component, and integration tests};
\end{tikzpicture}
\captionof{figure}{A testing pyramid for ML systems. Lower layers run more
often; higher layers validate consequential behaviour.}
\label{fig:ml-testing-pyramid}
\end{center}

\begin{table}[htbp]
\centering
\caption{Examples of tests at different ML system layers.}
\label{tab:ml-test-layers}
\begin{tblr}{
    width=\textwidth,
    colspec={Q[l,0.18\textwidth] X[l] X[l]},
    row{1}={font=\bfseries, bg=\maincolor!15!white},
    hlines,
    vlines}
Layer & Example assertion & Failure detected \\
Data contract & Schema, units, ranges, missingness, uniqueness &
Broken ingestion or silent source changes \\
Numerical & Analytical gradient agrees with finite differences &
Incorrect derivatives, broadcasting, or sign errors \\
Training sanity & A small model can overfit a tiny clean batch &
Disconnected loss, frozen parameters, or label mismatch \\
Behavioural & Invariance under irrelevant formatting or permutation &
Shortcut learning and brittle preprocessing \\
Robustness & Quality under corruption, shift, and adversarial cases &
Unsafe operating boundaries \\
Integration & Saved model reproduces accepted predictions after reload &
Serialization, preprocessing, or version mismatch \\
System & Latency, memory, concurrency, fallback, and rollback budgets &
Deployment failures invisible to offline accuracy \\
\end{tblr}
\end{table}

\subsection{Property and metamorphic tests}

Many ML outputs lack one exact expected value, but useful relationships remain
testable.
A unit conversion should change a physically meaningful output predictably.
Permuting a set should not change a permutation-invariant model.
Adding masked padding tokens should not alter unmasked sequence outputs.
Behavioural testing libraries formalise such capability-oriented checks for
NLP systems.\textsuperscript{\cite{ribeiroBeyondAccuracyBehavioral2020}}

\subsection{Protected regression gates}

Before deployment, define metrics that a new version must not degrade beyond a
declared tolerance.
Protect not only mean task quality but also safety tests, subgroup errors,
calibration, latency, memory, and cost.
The ML Test Score offers a production-oriented rubric spanning data,
features, model development, and
infrastructure.\textsuperscript{\cite{breckMLTestScore2017}}

\section{From Experiment to Operated System}

\keyterm{MLOps} applies software and operational discipline to the complete
machine-learning lifecycle.
Its goal is not merely automated deployment; it is controlled, traceable, and
reversible change.

\begin{center}
\begin{tikzpicture}[
    node distance=8mm and 8mm,
    stage/.style={draw=\maincolor!70!black,fill=\maincolor!6!white,
                  rounded corners=1mm,font=\sffamily\small,align=center,
                  minimum width=2.5cm,minimum height=9mm},
    >={Stealth[length=2mm]}]
    \node[stage] (scope) {Scope and\\risk};
    \node[stage,right=of scope] (data) {Versioned\\data};
    \node[stage,right=of data] (train) {Reproducible\\training};
    \node[stage,below=of train] (gate) {Evaluation\\gate};
    \node[stage,left=of gate] (deploy) {Controlled\\deployment};
    \node[stage,left=of deploy] (monitor) {Monitor and\\respond};
    \draw[->] (scope) -- (data);
    \draw[->] (data) -- (train);
    \draw[->] (train) -- (gate);
    \draw[->] (gate) -- (deploy);
    \draw[->] (deploy) -- (monitor);
    \draw[->] (monitor) -- (scope);
\end{tikzpicture}
\captionof{figure}{The controlled ML lifecycle: every release is measurable
and reversible.}
\label{fig:controlled-ml-lifecycle}
\end{center}

\subsection{What to version}

At minimum, record code commit, environment lockfile, data snapshot or query,
feature definitions, model configuration, seed, training logs, checkpoint,
evaluation predictions, and the command that reproduces the result.
A model file without its preprocessing and provenance is an archaeological
artifact, not a release.

\subsection{What to monitor}

Production monitoring covers several distinct questions:

\begin{itemize}
    \item \textbf{service health:} availability, errors, latency, throughput,
          memory, and cost;
    \item \textbf{input health:} schema violations, missingness, ranges,
          categories, and drift;
    \item \textbf{output health:} score distribution, abstentions, constraint
          violations, and unsafe content;
    \item \textbf{delayed quality:} task metrics when ground truth eventually
          arrives;
    \item \textbf{decision impact:} workload, user outcomes, feedback loops,
          and unintended behaviour.
\end{itemize}

Feature drift does not automatically imply quality loss, and stable features
do not prove stable quality.
Monitoring is a set of warning instruments, not a proof that the aircraft is
sound.

\section{Practical Lab: A Reproducible Model Comparison}

\begin{practicallab}[Leakage, baselines, uncertainty, and a release gate]
Choose a small tabular classification dataset with a meaningful group or time
variable.
Build one heuristic, logistic regression, a decision tree, and one ensemble.

\textbf{Required protocol:}
\begin{enumerate}
    \item write the evaluation contract and choose the split before fitting;
    \item put all learned preprocessing inside a pipeline;
    \item freeze a test set and tune on training/validation data only;
    \item report precision, recall, calibration, workload, and a paired
          bootstrap interval;
    \item introduce one deliberate leakage path and measure the false
          improvement;
    \item add data-contract, serialization, and behavioural tests;
    \item define a release gate including quality and latency;
    \item write a short model card with intended use, limitations, and evidence.
\end{enumerate}

\textbf{Exit criterion:}
another student can reproduce the comparison from one command and explain why
the selected model is preferable for the stated decision.
\end{practicallab}

\begin{checkpoint}
Why is a final test set not ``more data for tuning''?
Give one example where accuracy improves while the system becomes worse.
Name one invariant that your next ML project can test without knowing an exact
expected prediction.
\end{checkpoint}

\section{Further Reading}

The production testing rubric of Breck et al. connects model quality to system
readiness.\textsuperscript{\cite{breckMLTestScore2017}}
Sculley et al. explain why small modelling changes can create large systems
costs.\textsuperscript{\cite{sculleyHiddenTechnicalDebt2015}}
Datasheets and model cards provide complementary documentation structures for
data and trained
models.\textsuperscript{\cite{gebruDatasheetsDatasets2021,
mitchellModelCardsModel2019}}
